\documentclass[12pt]{article}
\usepackage[utf8]{inputenc}
\usepackage[T1]{fontenc}
\usepackage{hyperref}
\title{Tawala Language Resources for K-12 Curriculum Development}
\author{Compiled by Claude Sonnet 5 (claude-sonnet-5) for Andr\'as Kornai}
\date{}
\begin{document}
\maketitle

\section{Introduction}
Tawala (ISO 639-3 tbo; Glottocode tawa1275) is an Austronesian language of Milne Bay Province, Papua New Guinea, spoken on the East Cape peninsula, on the shores of Milne Bay and on parts of Sideia and Basilaki. Wikipedia gives about 20,000 speakers; Joshua Project gives 41,000. No school materials were found, only a reference grammar and indications of scripture.

\section{Primary Educational Resources}
None found. No dictionary was found in the search results either.

\section{Linguistic and Reference Materials}
\begin{itemize}
\item Ezard, Bryan (1997). \emph{A Grammar of Tawala}, Pacific Linguistics C-137 (free PDF, Australian National University). \url{https://openresearch-repository.anu.edu.au/bitstream/1885/145950/1/PL-C137.pdf}
\item Library record for the grammar (Papua New Guinea University of Technology). \url{https://matdisc.unitech.ac.pg/Record/74211}
\item Ezard 1997 source entry (TLOPO). \url{https://tlopo.clld.org/sources/Ezard1997}
\item Wikipedia, Tawala language. \url{https://en.wikipedia.org/wiki/Tawala_language}
\end{itemize}

\section{Religious and Community Resources}
Joshua Project lists a New Testament, an audio Bible and the Jesus Film for Tawala. No title or publisher was found, and the page was not read in detail.
\begin{itemize}
\item Joshua Project, Tawala language page. \url{https://joshuaproject.net/index.php/languages/tbo}
\end{itemize}

\section{Digital and Audio Resources}
Apart from the audio Bible listed by Joshua Project, no audio or digital resources were found.

\section{Curriculum Development Context}
No school programme or orthography document was found. The Ezard grammar is the main basis for writing a pedagogical description. Anyone searching for Tawala scripture should not confuse it with the Dawawa New Testament (\url{https://ebible.org/dww/FRT01.htm}), which is a different Milne Bay language.
\end{document}